\chapter{Verification Snapshot}

This appendix records the targeted verification work carried out while preparing the final cleaned repository and thesis workspace. The purpose is not to claim comprehensive software assurance. The purpose is to document that the most thesis-critical pathways and evidence surfaces were executed directly and observed to pass.

\input{tables/generated/verification_checks.tex}

\section{Training-Validity Evidence}

\begin{figure}[H]
    \centering
    \includegraphics[width=0.98\textwidth]{training_validity_evidence.pdf}
    \caption{Recovered W\&B evidence used to demonstrate training validity: optimization traces, evaluation traces, artifact availability, and raw-evidence inventory.}
    \label{fig:training-validity-evidence}
\end{figure}

\section{Interpretation of the Verification Checks}

The verification work now covers five risk classes. Canonical reporting rebuilds protect the thesis against stale or conflicting result summaries. Asset regeneration checks protect the link between the reporting layer and the LaTeX workspace. Full PDF rebuild checks protect the document pipeline itself. Count and stale-text scans protect the repository against exactly the sort of inconsistencies that a committee is likely to notice during review. The recovered W\&B evidence figure adds a direct training-validity surface that can be shown in a viva or committee discussion.

\section{Residual Verification Boundaries}

These checks still do not replace agronomic validation, and they do not prove that simulator performance differences would transfer to field performance. They are engineering and reporting checks. The remaining risk is therefore no longer that the campaign was unfinished, but that the completed evidence base still inherits the modeling boundaries described in Chapter~7, especially simulation-only evaluation, missing irrigation control, and the narrow crop-planning scope.
